% English translation of the Bulgarian source.
\documentclass[a4paper]{report}
% Shared preamble for the English probability and statistics documents.
% Based on the original course preamble. Compile with pdfLaTeX.
\usepackage[top=1.5cm,bottom=2cm,left=2cm,right=2cm]{geometry}
\usepackage[T1]{fontenc}
\usepackage{lmodern}
\usepackage[utf8]{inputenc}
\usepackage[english]{babel}

\usepackage{amsmath}
\usepackage{amssymb}
\usepackage{amsthm}
\usepackage{mathrsfs}
\usepackage{mathtools}
\usepackage{etoolbox}
\usepackage{tikz}
\usetikzlibrary{shapes.multipart,arrows.meta}
\usepackage{xcolor}
\usepackage{enumerate}
\usepackage[shortlabels]{enumitem}
\usepackage{hyperref}
\usepackage{version}
\usepackage{bbold}
\usepackage{fontawesome5}
\usepackage{fancyhdr}
\lfoot{}
\usepackage{natbib}
\usepackage{graphicx}
\usepackage{pdfpages}
\usepackage{authblk}
    
\usepackage{blindtext}
\usepackage{multicol}
\usepackage{collcell}
\usepackage{wrapfig}
\usepackage{blkarray}
\usepackage{cancel}
\usepackage{appendix}
\usepackage{float}
\makeatletter
\input{tex/assets/slashbox.sty}
\makeatother
\hypersetup{unicode,colorlinks,citecolor=red,filecolor=black,
  linkcolor=red,urlcolor=cyan}
\setlength{\columnsep}{1cm}
\pagenumbering{gobble}

\def\numtext#1{%
  \ifcase#1\or one\or two\or three\or four\or five\or six\or seven\or eight\or nine\or ten\or
  eleven\or twelve\or thirteen\or fourteen\or fifteen\or sixteen\or seventeen\or eighteen\or nineteen\or twenty\or Lots\fi}

\newcommand{\Byte}[1]{%
  \def\nodecontents{}%
  \foreach[count=\l] \k in {#1} {\xappto{\nodecontents}{\noexpand\nodepart{\numtext{\l}}\k}}%
  \node[rectangle split,rectangle split parts=8,rectangle split horizontal,draw] at (2,2){\nodecontents};}

\newcommand{\D}{{\mathrm{d}}}
\newcommand{\cE}{{\mathcal{E}}}
\newcommand{\cG}{{\mathcal{G}}}
\newcommand{\cX}{{\mathcal{X}}}
\newcommand{\cY}{{\mathcal{Y}}}
\newcommand{\cB}{{\mathcal{B}}}
\newcommand{\cF}{{\mathcal{F}}}
\newcommand{\cO}{{\mathcal{O}}}
\newcommand{\cS}{{\mathcal{S}}}
\newcommand{\cT}{{\mathcal{T}}}
\newcommand{\given}{\, \big| \,}
\newcommand{\cL}{{\mathscr{L}}}
\newcommand{\cD}{{\mathcal{D}}}
\newcommand{\cA}{{\mathscr{A}}}
\newcommand{\cP}{{\mathscr{P}}}
\newcommand{\cR}{{\mathscr{R}}}
\newcommand{\A}{{\mathcal A}}
\newcommand{\one}{\mathbb{1}}
\newcommand{\mathbbm}[1]{\mathbb{#1}}

\providecommand{\C}{}
\renewcommand{\C}{\mathbb C}
\newcommand{\R}{\mathbb R}
\newcommand{\Q}{\mathbb Q}
\newcommand{\N}{\mathbb N}
\newcommand{\T}{\mathbb T}
\newcommand{\F}{\mathbb F}
\renewcommand{\P}{\mathbb P}
\newcommand{\Z}{\mathbb Z}
\newcommand{\E}{\mathbb {E}}
\newcommand{\Pbb}[1]{\P\left( #1\right)}
\newcommand{\Ebb}[1]{\E\left[ #1\right]}

\newcommand{\lb}{\left (}
\newcommand{\rb}{\right )}
\newcommand{\lbb}{\left [}
\newcommand{\rbb}{\right ]}
\newcommand{\labs}{\left |}
\newcommand{\rabs}{\right |}
\newcommand{\lbrb}[1]{\lb #1 \rb}
\newcommand{\lbbrbb}[1]{\lbb#1\rbb}
\newcommand{\labsrabs}[1]{\labs#1\rabs}
\newcommand{\lbbrb}[1]{\lbb#1\rb}
\newcommand{\lbrbb}[1]{\lb#1\rbb}
\newcommand{\langlerangle}[1]{\langle#1\rangle}
\newcommand{\lbcurly}{\{}
\newcommand{\rbcurly}{\}}
\newcommand{\curly}[1]{\lbcurly#1\rbcurly}
\newcommand{\abs}{\labsrabs}
\newcommand{\ind}[1]{1_{#1}}
\newcommand{\indep}{\perp \!\!\! \perp}
\renewcommand{\d}{\, \mathrm{d}}

\newtheorem{theorem}{Theorem}
\newtheorem{problemCnt}{Lemma}
\theoremstyle{definition}
\newtheorem{problem}[problemCnt]{Problem}
\newtheorem{problembook}[problemCnt]{\faBook \,  Problem}
\newtheorem{problempencil}[problemCnt]{\faPencil* \,  Problem}
\newtheorem{problemstar}[problemCnt]{\faStar\,  Problem}
\newtheorem{problemstarpencil}[problemCnt]{\faStar \, \faPencil* \,   Problem}
\newtheorem{probleminteresting}[problemCnt]{\faSmile[regular] \,   Problem}
\newtheorem{solution}{Solution}

\newcommand{\fwcell}[1]{\makebox[\arraystretch\normalbaselineskip]{$#1$}}
\newenvironment{enumeratei}{\begin{enumerate}[(i)]}{\end{enumerate}}
\newenvironment{enumeratea}{\begin{enumerate}[(a)]}{\end{enumerate}}
\newenvironment{enumeraten}{\begin{enumerate}[\upshape 1.]}{\end{enumerate}}

% Compatibility helpers needed by a few files in the combined archive.
\AtBeginDocument{\noindent{\footnotesize English translation prepared with ChatGPT (OpenAI, GPT-6), 26 October 2026, for translation of the Bulgarian text into English and LaTeX formatting.
The Bulgarian originals remain the reference; translation may contain errors.}\par\smallskip}
\newcommand{\MissingGraphic}[1]{%
  \begin{center}\fbox{\parbox{0.82\linewidth}{\centering Missing source image: \texttt{#1}}}\end{center}}



\begin{document}
\noindent Probability and Statistics / SEM, 24 August 2025 \hfill Examination problems
\hrulefill
\begin{it}
The exact grading formula will depend on the results. As an approximation, use 2 plus the number of points. Time allowed: 3 hours. Good luck.
You may use $\sum_{k\geq1}\frac{x^k}{k}=-\ln(1-x)$ for $|x|<1$.
\end{it}
\begin{problem}
\begin{enumerate}
\item (0.25 points) An urn contains $n$ balls, and all possible numbers of white balls are equally likely. We add one white ball to the urn. What is the probability that a randomly selected ball from the urn is white?
\item (0.25 points) An urn contains $n$ white and $m$ black balls. Two players draw alternately from the urn, replacing each drawn ball. The first to draw a white ball wins. What is the probability that the first player wins? What happens as $m\to\infty$?
\item (0.25 points) We have two urns, 5 white balls and 5 black balls. How should we distribute the balls by number and colour between the urns to maximize the probability of drawing a white ball when we choose a random urn and then a random ball from it?
\item (0.5 points) Ivan has only 10-stotinka coins. The price of every item in a shop, in stotinki, ends in a random digit from $\{0,1,2,\dots,9\}$, with all final digits equally likely and independent. Ivan buys $n\geq1$ items. What is the probability that he can pay exactly, without receiving change\footnote{That is, his bill in stotinki is a multiple of 10.}? What is the expected number of items he must take until he can pay exactly for the first time\footnote{That is, the expected number of items until the total price in stotinki first becomes a multiple of 10.}?
\end{enumerate}
\end{problem}
\begin{problem}
Consider a Bernoulli scheme $(n,p)$. Let $\mu_n$ be the number of successful trials. The number $\tilde p_n:=\mu_n/n$ is called the relative frequency of successes.
\begin{enumerate}
\item (0.25 points) If $p=0.4$ and $n=1500$, what is the probability that $\tilde p_n$ lies in $(0.4,0.44)$?
\item (0.25 points) If $p=0.375$, how many independent trials are needed so that $\{|\tilde p_n-p|\leq0.01\}$ has probability at least $0.995$?
\item (0.25 points) Let $p=2/3$ and $n=800$. What $\varepsilon$ should we choose so that $\{|\tilde p_n-p|<\varepsilon\}$ has probability at least $0.985$?
\item (0.25 points) If $n=10000$, for what value of $p$ will the probability of $\{|\tilde p_n-p|\leq0.01\}$ be at least $0.99$?
\end{enumerate}
\end{problem}
\begin{problem}
Let $p_k=ba^k/k$ for $k\geq1$ and $0<a<1$.
\begin{enumerate}
\item (0.25 points) For what value of $b$ is $(p_k)_{k\geq1}$ a probability distribution?
\item (0.25 points) Find its probability generating function, expectation and variance.
\item (0.25 points) Let $X$ have the distribution above. Find the probability generating function of $2X+1$.
\end{enumerate}
\end{problem}
\begin{problem}
\begin{enumerate}
\item (0.75 points) Three random points are chosen independently and uniformly on a circle of radius 1. Calculate the probability that the length $L$ of the longest side of the triangle with these vertices does not exceed $\sqrt3$.
\item (0.25 points) Consider a generalization: how can we calculate $\Pbb{L<\varepsilon}$, where $\varepsilon\leq\sqrt3$?
\end{enumerate}
\end{problem}
\begin{problem}
Let $U\sim U(0,1)$ and $V\sim U(-\pi/2,\pi/2)$.
\begin{enumerate}
\item (0.25 points) One way to generate a pseudorandom number between 0 and 1 using $U$ is to choose a natural number $M$, usually large, and calculate the fractional part of $MU$: $X:=\{MU\}$. Find $\E X$, $DX$ and $Cor(U,X)$. What do you think of this method?
\item (0.25 points) Let $U_1,\dots,U_{100}\sim U$ be iid. What is the expectation of each? What about the largest and smallest among them?
\item (0.25 points) Find the density, expectation and variance of $Y:=\tan V$.
\item (0.25 points) Let $Z\sim Cauchy(1)$, that is, $f_Z(x)=1/(\pi(1+x^2))$ for $x\in\R$. Prove that the random variables $1/Z$, $2Z/(1-Z^2)$ and $(3Z-Z^3)/(1-3Z^2)$ have the same distribution.
\end{enumerate}
\end{problem}
\end{document}
